% Capitulo: Modelagem para ciência de dados
\chapter{Modelagem para ciência de dados}
\label{cap:modelagem-cd}

O esquema estrela é ótimo para somar, filtrar e comparar. Para treinar um modelo que prevê o
futuro, ele quase nunca serve do jeito que está. Modelos preditivos precisam de outra forma
de dado, com outro grão e outros cuidados, e modelos prescritivos precisam ainda de um tipo
de dado que a análise descritiva nunca guardou: o registro das decisões tomadas. Este
capítulo trata dessas diferenças e do erro mais caro da modelagem para ciência de dados, o
vazamento de informação do futuro para o passado.

\section{Descritiva, preditiva e prescritiva}
\label{sec:tres-niveis}

É comum separar a análise de dados em três níveis de ambição. A descritiva responde ao que
aconteceu. A preditiva responde ao que vai acontecer. A prescritiva responde ao que fazer a
respeito. Uma das formulações mais conhecidas dessa escada, do relatório padrão à
otimização, vem de \textcite{davenport2007}, que associam cada degrau a uma vantagem
competitiva maior para quem chega lá.

Um exemplo de logística ajuda a ver a diferença. A análise descritiva mostra que, no último
trimestre, as entregas de uma certa região atrasaram mais do que as das outras, e
principalmente às segundas-feiras. A preditiva estima, para cada entrega programada para
amanhã, a probabilidade de ela atrasar. A prescritiva usa essas probabilidades, junto com o
custo de cada veículo, a capacidade disponível e as janelas de cada cliente, para recomendar
como montar as rotas de forma a minimizar atrasos e custo ao mesmo tempo.

Cada nível depende do anterior. Não se prevê bem o que não se descreve bem: se o atraso é
registrado de formas diferentes em cada filial, o modelo aprende a diferença entre as
filiais, e não o atraso. E não se prescreve bem sem uma previsão confiável. Por isso a
qualidade da camada descritiva, com tudo o que os capítulos anteriores discutiram, é
pré-requisito para os outros dois.

\rascunho{
  Opcional: ligar com a minha experiência de pesquisa, que segue os três níveis
  (descritivo, preditivo, prescritivo), sem detalhar dados ou empresa.
}

\section{A tabela de variáveis}
\label{sec:tabela-features}

Um modelo preditivo aprende a partir de exemplos. Cada exemplo é uma linha com as
informações disponíveis num certo momento, as variáveis (\emph{features}), e o que aconteceu
depois, o alvo. Para prever se uma entrega vai atrasar, cada linha seria uma entrega, com
variáveis como a região, o dia da semana, a distância, o histórico recente de atrasos do
motorista e da rota, e o alvo: atrasou ou não.

Isso muda o grão. No modelo dimensional, a tabela de fatos tem uma linha por evento (cada
entrega realizada), e as perguntas são respondidas agregando. Na tabela de variáveis, o grão
é \enquote{uma linha por entidade e por momento de previsão}: um cliente por mês, um veículo
por dia, uma entrega no momento em que foi programada. E as variáveis muitas vezes já são
agregações calculadas para aquele momento, como o número de atrasos da rota nos trinta dias
anteriores.

O resultado são tabelas largas, com dezenas ou centenas de colunas, que um modelador
dimensional evitaria. Não há problema nisso: a tabela de variáveis é um produto da camada
ouro feito para um consumidor específico, o modelo, e deve ser desenhada para ele. Na
arquitetura medalhão, ela costuma ser construída a partir da prata, e não do ouro do BI,
porque o ouro já vem agregado e perdeu o detalhe de que o modelo precisa
(\veja{sec:bronze-prata-ouro}).

\section{O futuro não pode vazar para o passado}
\label{sec:vazamento}

O erro mais comum, e mais caro, na construção de tabelas de variáveis é o vazamento
(\emph{data leakage}): usar, para treinar o modelo, uma informação que não estaria disponível
no momento em que a previsão precisa ser feita \parencite{kaufman2012}. O modelo fica
excelente nos testes, porque está, na prática, olhando a resposta. Em produção, onde essa
informação não existe ainda, ele erra.

O vazamento raramente é óbvio. Uma variável como \enquote{número de reclamações do cliente}
parece inocente, mas, se calculada com todos os dados disponíveis hoje, inclui as
reclamações feitas depois da data que se quer prever, inclusive aquelas causadas justamente
pelo atraso que o modelo deveria antecipar. Um campo de status que é atualizado quando a
entrega termina vaza o alvo diretamente. E até a forma de separar treino e teste pode vazar,
se dados do futuro entram no treino e dados do passado no teste \parencite{huyen2022}.

A defesa é a correção no tempo (\emph{point-in-time}): cada variável é calculada só com o
que se sabia na data de referência do exemplo. A consulta abaixo mostra a ideia, calculando
para cada cliente e cada data de referência as compras dos noventa dias anteriores, e nada
depois.

\begin{sql}
-- Variáveis de cada cliente calculadas só com o passado de cada data de referência
SELECT
    r.id_cliente,
    r.data_referencia,
    COUNT(v.id_venda)                AS compras_90d,
    COALESCE(SUM(v.valor_total), 0)  AS valor_90d,
    r.cancelou_no_mes_seguinte       AS alvo
FROM referencias_clientes AS r
LEFT JOIN vendas AS v
       ON v.id_cliente = r.id_cliente
      AND v.data_venda <  r.data_referencia                       -- nada do futuro
      AND v.data_venda >= r.data_referencia - INTERVAL '90' DAY
GROUP BY r.id_cliente, r.data_referencia, r.cancelou_no_mes_seguinte;
\end{sql}

A correção no tempo também vale para os atributos. Se o cliente mudou de segmento no meio do
ano, o modelo precisa saber o segmento que ele tinha na data de referência, e não o atual.
Aqui reaparece uma técnica do capítulo~\ref{cap:dw}: as dimensões que mudam devagar, com
uma linha por versão e datas de validade (\veja{sec:kimball}). O que no BI é um recurso para
análise histórica, na ciência de dados vira requisito.

\begin{destaque}
Se o modelo parece bom demais nos testes, a primeira hipótese não é que ele é ótimo. É que
algum dado do futuro vazou para o treino.
\end{destaque}

\section{Reuso de variáveis e dívida técnica}
\label{sec:feature-store}

Quando cada cientista de dados calcula as próprias variáveis, o mesmo conceito ganha várias
definições. \enquote{Cliente ativo} é calculado de um jeito no modelo de cancelamento e de
outro no de recomendação, e ninguém sabe qual está certo. É o mesmo problema das métricas
divergentes que o data warehouse tentou resolver nos anos 1990, agora do lado da ciência
de dados.

A resposta também é parecida: definir uma vez e reaproveitar. Os repositórios de variáveis
(\emph{feature stores}) guardam as definições, calculam as variáveis com correção no tempo
para o treino e servem os mesmos valores para o modelo em produção, garantindo que ele veja
na operação exatamente o que viu no treino \parencite{huyen2022}. Nem todo time precisa de
uma ferramenta dedicada para isso; uma camada bem organizada de tabelas de variáveis, com
testes e documentação, resolve muitos casos.

Vale lembrar uma observação que ficou famosa entre quem constrói sistemas de aprendizado de
máquina: o código do modelo propriamente dito é uma parte pequena do sistema. Em volta dele
há coleta, preparação, validação e monitoramento de dados, e é nessa parte maior que a dívida
técnica se acumula \parencite{sculley2015}. Para quem vem da engenharia e da análise de
dados, é uma boa notícia. A maior parte do trabalho de colocar um modelo em produção é
trabalho de dados, e as práticas da Parte~\ref{parte:devops} valem integralmente.

\section{Da previsão à decisão}
\label{sec:prescritiva}

A análise prescritiva junta a previsão com o que o negócio pode fazer. Isso exige dados que
as camadas anteriores raramente têm: os custos de cada alternativa, as capacidades (quantos
veículos, quantas pessoas, quanto estoque), as restrições (janelas de entrega, contratos,
regras trabalhistas) e os objetivos, muitas vezes em conflito, como custo e nível de
serviço. Com isso, a recomendação costuma sair de um problema de otimização, e não só de um
modelo estatístico.

Há um dado novo que a prescritiva exige e que quase nenhuma empresa guarda: o registro das
decisões. Para saber se a recomendação funciona, é preciso registrar o que foi recomendado,
o que de fato foi decidido (as pessoas nem sempre seguem a recomendação, e muitas vezes têm
bons motivos) e o que aconteceu depois. Sem esse registro, não há como avaliar o sistema nem
melhorá-lo. Com ele, cada decisão vira um exemplo para a próxima.

Vale também modelar cenários como dados, e não como planilhas soltas. Quando o gestor
pergunta \enquote{e se a demanda crescer 20\%?}, a resposta deveria vir de um conjunto de
premissas versionado, comparável com outros cenários, e não de uma cópia da planilha com um
número alterado.

Por fim, as restrições certas vêm de quem conhece a operação. Um modelo prescritivo que
ignora que um cliente não recebe entregas depois das 16 horas, ou que um motorista não pode
dirigir mais de um certo número de horas seguidas, recomenda o impossível. É mais um lugar
em que o conhecimento de negócio (capítulo~\ref{cap:negocio}) é o que separa o modelo
elegante do modelo útil.

\begin{leituras}
  \leitura{huyen2022}{A visão mais completa e prática de como dados, variáveis e
    modelos se organizam num sistema de aprendizado de máquina em produção.}
  \leitura{kaufman2012}{O artigo de referência sobre vazamento de dados: como ele acontece,
    como detectar e como evitar.}
  \leitura{sculley2015}{Artigo curto que mostra por que o modelo é a menor parte do
    problema e onde a dívida técnica se esconde.}
\end{leituras}
